\chapter{Impostazione sperimentale}
\label{sec:setup}

% Source notes: Project/docs/thesis_notes.md; checkpoint_2026-10-10.md; configs Project/GroundOne/configs/train_0.5b.yaml,
% train_3b.yaml; queues Project/jobs/queue_night3.txt, queue_day4.txt, queue_night5.txt, queue_day5.txt, queue_night6.txt.

Questo capitolo descrive i modelli addestrati, i sistemi con cui vengono confrontati, le metriche, il protocollo
statistico e il modo in cui si misurano tempo ed energia. Tutti gli esperimenti sono eseguiti su un'unica macchina, un
MacBook Pro con chip Apple M4 (CPU a 10 core, GPU a 10 core, 120\,GB/s di banda di memoria) e 16\,GB di memoria
unificata~\cite{apple2024m4specs}: in questo modo ogni risultato riflette ciò che un verificatore è in grado di fare sul
dispositivo su cui gira un assistente locale. Anche ogni baseline, compresi gli encoder e il System One generalista, è
eseguita su questo portatile; nessun risultato dipende da un servizio remoto.

\section{Modelli}
\label{sec:setup-models}

Entrambi i modelli GroundOne partono dai checkpoint di Qwen2.5 addestrati a seguire istruzioni~\cite{qwen2024qwen25},
convertiti in MLX con pesi a 8 bit, ed entrambi sono addestrati con la procedura della
Sezione~\ref{sec:method-training}. Differiscono soltanto per dimensione, per numero di risposte di addestramento e per il
modo in cui si formano i micro-batch, aspetto dettato dalla memoria disponibile.

Il modello 0.5B è addestrato su tutte le 13{,}203 risposte di addestramento. Un'epoca di 802 passi dell'ottimizzatore ha
richiesto 5.5 ore di tempo reale (wall-clock), con una valutazione sull'intero insieme di sviluppo ogni 150 passi, e
l'addestramento ha occupato al massimo 2.6\,GB di memoria. Il modello 3B è circa sei volte più lento per token, perciò
viene addestrato per una notte su un sottoinsieme casuale di 6{,}500 risposte (circa 45{,}000 decisioni): 403 passi in
12.25 ore, con un picco di 5.9\,GB. La sua valutazione sullo sviluppo, ogni 80 passi, usa un sottoinsieme fisso di 200
risposte, poiché un passaggio completo richiederebbe circa un'ora. In entrambi i casi si conserva il checkpoint con il
miglior punteggio sullo sviluppo, che viene poi calibrato come descritto nella Sezione~\ref{sec:method-calibration}; per
il 3B si tratta del passo 320.

Per misurare quanto i risultati dipendano dal seed casuale, che determina l'inizializzazione degli adapter e l'ordine dei
batch, la procedura del 0.5B viene ripetuta con un secondo seed. È noto, infatti, che gli esiti del fine-tuning variano
con entrambi i fattori~\cite{dodge2020finetuning}. Un'ulteriore esecuzione con un terzo seed è andata precocemente in
collasso verso una predizione costante ed è stata interrotta a circa il 60\,\% dell'epoca; essendo incompleta, è esclusa
dalle tabelle e viene discussa nella Sezione~\ref{sec:results-seeds}. La seconda esecuzione, portata a termine, è valutata
su tutti e tre gli insiemi di test esattamente come la prima.

\section{Baseline}
\label{sec:setup-baselines}

Le baseline servono a isolare le fonti di un eventuale vantaggio, e ciascuna risponde a una domanda specifica.

\paragraph{Indizi che non leggono mai il documento nel suo insieme.} Due predittori semplici stabiliscono la soglia
minima che qualunque verificatore deve superare. Il primo assegna a una frase un punteggio pari alla quota delle sue
parole di contenuto che compaiono anche nel documento. Il secondo è un classificatore bag-of-words addestrato soltanto
sulle frasi di addestramento, senza il documento. Entrambi sono descritti nella Sezione~\ref{sec:data-eda}.

\paragraph{Readout senza addestramento (B1).} Agli stessi modelli Qwen2.5, privi di fine-tuning, viene posta la domanda a
tre classi per ciascuna frase e se ne leggono i logit delle lettere di risposta (readout, ossia lettura dei logit).
Poiché un modello non addestrato non ha motivo di interpretare il marcatore \texttt{Verdict:} come una domanda, questa
baseline adotta il formato con una domanda esplicita per frase (Sezione~\ref{sec:method-readout}). Essa misura ciò che il
meccanismo System One offre prima di qualsiasi specializzazione.

\paragraph{Giudici generativi (B0a, B0b).} I modelli non addestrati sono impiegati anche nel modo più comune nella
pratica: ricevono il documento, la risposta numerata e l'elenco delle frasi da controllare, e scrivono un verdetto per
ciascuna frase con decodifica greedy. Il giudice \emph{sintetico} scrive una riga ``\texttt{<number>: <letter>}'' per
frase, mentre il giudice \emph{verboso} scrive una breve frase di ragionamento seguita dal verdetto. Un verdetto che non
può essere interpretato costituisce un errore di formato; viene conteggiato come errore nell'accuratezza e riportato
separatamente.

\paragraph{Giudici a parità di output (B0d).} I giudici sintetico e verboso restituiscono meno di GroundOne: una lettera
per frase e nessuna probabilità, e quindi non offrono calibrazione, soglia regolabile né scelta della politica. A due
ulteriori giudici si chiede pertanto esattamente il contenuto prodotto da GroundOne, cioè un'etichetta e tre probabilità
per frase. Il primo li scrive come oggetto JSON, con le probabilità espresse a parole, che è ciò che uno sviluppatore
otterrebbe con il solo prompt. Il secondo è il giudice sintetico le cui tre probabilità vengono lette dai logit in
corrispondenza della lettera generata; ottiene lo stesso contenuto al costo del giudice sintetico ed è l'avversario
generativo più forte. Entrambi i giudici sono calibrati con la stessa temperatura e lo stesso bias di GroundOne, stimati
su 200 risposte di sviluppo, affinché le loro probabilità ricevano lo stesso trattamento. Il prompt JSON è stato scelto su
dodici risposte di sviluppo unicamente in base alla validità del formato: un esempio numerico presente nel prompt veniva
copiato alla lettera dai modelli ed è stato quindi sostituito da segnaposto. L'interpretazione dell'output è permissiva, a
favore dei giudici: le voci complete di un oggetto troncato vengono conservate e un'etichetta mancante viene ricavata
dalle probabilità. La decodifica non è vincolata da una grammatica~\cite{willard2023outlines}, così che il protocollo
resti quello degli altri giudici; questa scelta incide soprattutto sugli errori di formato del modello più piccolo.

\paragraph{Readout contro generazione (B2).} I modelli GroundOne addestrati sono eseguiti con lo stesso prompt e lo
stesso singolo passaggio, ma il verdetto in corrispondenza di ciascun marcatore è il token con la probabilità più alta
sull'intero vocabolario, ossia ciò che produrrebbe la generazione greedy di un token. Le tre lettere, con o senza spazio
iniziale, corrispondono alle tre classi, e qualsiasi altro token è un errore di formato. Questo controllo mantiene
identici pesi, prompt e costo, e cambia soltanto il fatto che la decisione sia o meno ristretta alle opzioni ammissibili.

\paragraph{Un System One generalista (B3).} Il progetto aperto Rizzo-Flow~\cite{rizzoflow} realizza l'idea di System One
con un modello da 4B addestrato su decisioni tipizzate generiche. Gira in un proprio server, basato su
llama.cpp~\cite{llamacpp}, e viene interrogato via HTTP con il documento e la risposta numerata come stato condiviso e una
domanda a scelta a tre classi per frase. Si impiega un modello quantizzato a 4 bit perché quello a 8 bit non trova posto
in memoria accanto al sistema operativo.

\paragraph{Un modello decisionale generalista della stessa dimensione (B3b).} d1-3B di Liquid AI~\cite{liquid2026d1} è
un modello decisionale aperto da 3.1B parametri, post-addestrato per decisioni tipizzate su un'ampia miscela di compiti,
che legge lettere di risposta ristrette con lo stato condiviso tra le domande: appartiene alla stessa famiglia di readout
di GroundOne e ha la dimensione del suo modello più grande, ma è generalista. Permette quindi di isolare il valore della
specializzazione rispetto a un modello decisionale già addestrato a decidere, anziché rispetto a uno non addestrato.
Viene eseguito zero-shot con il proprio codice di inferenza, in PyTorch sulla GPU del portatile e senza accesso alla rete,
con lo stesso stato e le stesse domande di Rizzo-Flow e al massimo sei domande per chiamata sullo stesso stato; su
LLM-AggreFact i suoi documenti sono suddivisi esattamente come quelli di GroundOne e si conserva il massimo tra i
frammenti. Un dettaglio pratico è rilevante per la riproducibilità: PyTorch conserva kernel compilati per ogni nuova
lunghezza di input sulla GPU Apple e non rilascia quella memoria, perciò la valutazione si svolge in processi worker che
vengono riavviati quando la memoria usata dal driver della GPU supera 11\,GB, il che lascia invariati gli output.

\paragraph{Encoder specializzati (B4).} MiniCheck-RoBERTa~\cite{tang2024minicheck} e HHEM-2.1-open~\cite{vectara2024hhem}
sono i verificatori di piccole dimensioni più forti disponibili. Valutano ciascuna coppia (documento, frase) in modo
indipendente e restituiscono una sola probabilità di supporto; i documenti lunghi sono suddivisi in frammenti di circa 400
parole e si conserva il massimo tra i frammenti, come in MiniCheck. Non disponendo di un output a tre classi, sono
confrontati soltanto sulla decisione binaria. Girano con PyTorch sulla GPU del portatile.

\paragraph{Risultati pubblicati.} Per LLM-AggreFact, i punteggi di sistemi più grandi tratti dalla classifica pubblica,
come Qwen2.5-7B usato zero-shot, sono riportati così come pubblicati e non vengono rieseguiti.

I sistemi più lenti sono valutati su sottocampioni stratificati degli insiemi di test: Rizzo-Flow su 300 risposte di T1,
600 coppie di T3 e 1{,}100 elementi di T2, di cui ne ha valutati 1{,}093 (sette superavano il suo contesto); d1-3B su un
sottocampione casuale di 4{,}000 elementi di T2; il 3B non addestrato su 600 risposte di T1; i giudici generativi, infine,
sulle 200 risposte del protocollo di latenza descritto più avanti. GroundOne, il 0.5B non addestrato, gli encoder e, su T1
e T3, d1-3B sono valutati sugli insiemi di test completi, e i confronti con un sistema sottocampionato sono sempre
condotti sugli elementi valutati da entrambi i sistemi.

\section{Metriche}
\label{sec:setup-metrics}

Poiché la quota di frasi allucinate differisce tra l'insieme di sviluppo e T1 (Sezione~\ref{sec:data-eda}), le metriche
primarie non dipendono dal tasso di base. Per la decisione a tre classi la metrica principale è l'accuratezza bilanciata,
cioè la recall media sulle tre classi, affiancata dalla recall di ciascuna classe e dalla quota di ciascuna classe
predetta, che rivela un modello in collasso su un'unica risposta. Per la decisione binaria le metriche sono l'accuratezza
bilanciata con soglia 0.5, come nella classifica di LLM-AggreFact, e l'AUROC di $P(\text{supported})$, che non dipende da
alcuna soglia. La calibrazione è misurata con l'ECE definito nella Sezione~\ref{sec:method-calibration}, con il Brier
score e con la log-verosimiglianza negativa, mentre la predizione selettiva è valutata mediante curve rischio--copertura e
l'area sottesa a tali curve (AURC)~\cite{geifman2017selective}. Tutte le metriche su T1 sono riportate anche per compito e
per fascia di $K$.

\section{Protocollo statistico}
\label{sec:setup-statistics}

Ogni differenza tra due sistemi è accompagnata da un intervallo di confidenza al 95\,\% ottenuto con un bootstrap
appaiato con 2{,}000 ricampionamenti~\cite{koehn2004significance,bergkirkpatrick2012significance,dror2018hitchhiker}.
L'unità di ricampionamento non è la frase ma il gruppo usato dal controllo di leakage: il documento sorgente in RAGTruth,
il documento in LLM-AggreFact e la pagina di Wikipedia in VitaminC. Le frasi di una stessa risposta, così come le risposte
relative a uno stesso documento, sono fortemente correlate, e ricampionarle in modo indipendente produrrebbe intervalli
troppo stretti. Su LLM-AggreFact i gruppi sono ricampionati all'interno di ciascuno degli undici dataset, perché la media
macro attribuisce a ogni dataset lo stesso peso. Entrambi i sistemi di un confronto sono valutati sugli stessi
ricampionamenti, e il $p$-value riportato è la quota bilaterale di ricampionamenti in cui la differenza cambia segno.

Due limiti di questo protocollo vengono dichiarati in anticipo. In primo luogo, si sottopongono a test circa cinquanta
confronti senza alcuna correzione per confronti multipli; le conclusioni poggiano quindi su intervalli ben lontani dallo
zero, e non si rivendica alcuna differenza con $p$-value vicino a 0.05. In secondo luogo, il ricampionamento dei dati non
cattura la variabilità dovuta all'addestramento: questa informazione è fornita separatamente dalle due esecuzioni del
0.5B (Sezione~\ref{sec:results-seeds}).

\section{Latenza, energia e hardware}
\label{sec:setup-protocol}

La latenza è misurata su 200 risposte di T1, 50 per ciascuna fascia di $K$, così che le risposte lunghe siano
rappresentate quanto quelle brevi. GroundOne, il readout con una domanda per frase e tutti i giudici generativi girano nel
processo stesso su MLX, una risposta alla volta, come farebbe un verificatore all'interno di un assistente locale. Le
prime tre risposte servono da riscaldamento (warm-up) e sono escluse. Per le varianti System One ogni risposta è
cronometrata tre volte e si conserva la mediana; i giudici generativi sono cronometrati una sola volta, poiché la
decodifica greedy è deterministica e ogni esecuzione è molto più lunga. Il tempo riportato è il tempo reale (wall-clock)
della chiamata al modello per una risposta, da cui si ricava il tempo per decisione.

Sono stati registrati tempi anche per gli encoder, per Rizzo-Flow e per d1-3B, ma in condizioni che non consentono un
confronto equo. Gli encoder giravano con PyTorch e valutavano ogni frase separatamente, rileggendo ogni volta il
documento, il che li penalizza; Rizzo-Flow girava in un motore di inferenza diverso, con pesi a 4 bit e dietro
un'interfaccia HTTP; d1-3B girava in PyTorch con il proprio codice e con i riavvii dei worker descritti sopra. Non si
avanza quindi alcuna affermazione sulla latenza nei loro confronti.

L'energia è misurata con lo strumento \texttt{powermetrics} di macOS, che riporta la potenza di CPU, GPU e Neural Engine
ogni 200 millisecondi. Per ciascun sistema, 100 risposte del sottoinsieme di latenza (25 per fascia di $K$) sono
verificate una sola volta, dopo 60 secondi di misura a riposo con il modello già caricato; l'energia di una verifica è la
potenza media durante l'esecuzione, al netto della potenza a riposo, moltiplicata per la durata dell'esecuzione e divisa
per il numero di risposte. La misura copre le unità di calcolo del chip ma non la memoria, lo schermo o l'archiviazione,
ed è stata effettuata di notte con la macchina altrimenti inattiva.

Un solo modello è in esecuzione alla volta, e una guardia di memoria sorveglia ogni job lungo registrando per tutta la
sua durata la memoria libera e lo swap. I risultati vengono scritti una sola volta e mai sovrascritti, e ogni esecuzione
conserva la propria configurazione, così che ogni numero del Capitolo~\ref{sec:results} possa essere ricondotto
all'esecuzione che lo ha prodotto.
